\section{Translation and AST Preservation}
\label{app:translation}

\paragraph{Well-defined loop semantics.}
The transformer domain $\mathcal{H}$ from \Cref{sec:prelims} has pointwise suprema of increasing $\omega$-chains: these preserve $\omega$-continuity and the mass bound by monotone convergence. Its least element is $\bot=\lambda\eta.0$. By structural induction, the body transformer $\post{B}$ is $\omega$-continuous and does not increase mass. For $f\in\mathcal{H}$ and $\eta\in\subdists$,
\[
 \weight{\Phi_{b,B}(f)(\eta)}
 =\weight{[\neg b]\eta}
   +\weight{f\bigl(\post{B}([b]\eta)\bigr)}
 \leq\weight{[\neg b]\eta}+\weight{[b]\eta}
 =\weight{\eta}.
\]
Composition and addition preserve $\omega$-continuity here, so $\Phi_{b,B}(f)\in\mathcal{H}$. Moreover, for every increasing chain $(f_n)_n$ in $\mathcal{H}$, pointwise evaluation gives $\Phi_{b,B}(\sup_n f_n)=\sup_n\Phi_{b,B}(f_n)$. Kleene's theorem therefore yields \Cref{eq:while-transformer-fixpoint}. Structural induction also establishes linearity on sums that are subdistributions: in the loop case, every approximant is linear, and its pointwise supremum remains linear. The countable-sum property follows from $\omega$-continuity.

\paragraph{Operational correspondence.}
The translation $\mathcal{T}$ assigns distinct entry and exit locations to each command, with fresh internal locations for its subcommands. Skip is a single identity edge. Assignments are single update edges. A conditional entry has guarded edges to the two branch entries; a probabilistic-choice entry instead has edges of probabilities $p$ and $1-p$. In either case the branch exits are identified with the command exit. Sequential composition identifies the first command's exit with the second command's entry. For a while-loop, its entry is the guard: the true edge enters the body, the false edge reaches the program exit, and the body exit is identified with the guard. Zero-probability edges are ignored, and the program exit is absorbing. Thus the loop-free body is acyclic and every path through it contains at least one edge.

\begin{lemma}[Loop-free correspondence]
For loop-free $B$ and states $\sigma,\sigma'$,
\[
 \post{B}(\dirac{\sigma})(\sigma')
 =\Pr(\mathsf{Paths}_{\mathcal{T}[B]}((\ell_{\rm in},\sigma),(\ell_{\rm out},\sigma'))).
\]
\end{lemma}
\begin{proof}
By structural induction on $B$. Skip and assignment have one path. Sequential composition sums over intermediate states. Conditional paths are partitioned by the guard. Probabilistic choice multiplies branch-path probabilities by $p$ and $1-p$, matching linearity of the post-transformer.
\end{proof}

For a loop $C=\while{b}{B}$ with loop-free body, induction on $j$ shows that $\Phi_{b,B}^{j}(\bot)(\mu)$ is the subdistribution of pCFG runs that exit within at most $j$ visits to the outer loop guard. The zeroth approximant is zero; for $j\geq1$, this means termination after at most $j-1$ completed body iterations. The induction step partitions runs into those that exit at the first guard test and those that execute $B$ and then terminate within the remaining guard visits. Taking pointwise suprema gives equality between post-distribution weight and pCFG reachability probability. Therefore $\weight{\post{C}(\dirac{\sigma})}=1$ iff the translated graph reaches its exit almost surely from the loop guard with valuation $\sigma$. For any initial subdistribution,
\[
 \weight{\mu}-\weight{\post{C}(\mu)}
 =\sum_{\sigma\in\supp(\mu)}\mu(\sigma)
     \bigl(1-\weight{\post{C}(\dirac{\sigma})}\bigr).
\]
Every summand is nonnegative, so the left-hand side vanishes exactly when
$C$ is AST from every state in $\supp(\mu)$. Taking the union over
$\mu\in M$ proves the claimed equivalence for arbitrary $M$.

\subsection{Common Certificates for Initial-State Sets}
\label{app:initial-sets}

The pCFG completeness construction
of~\cite{DBLP:journals/pacmpl/MajumdarS25} is presented for a fixed initial
configuration. We record why it also supplies a single pair of
certificates for the initial-state sets used here. We use the purely probabilistic model with fixed rational probabilities.
Computability of the operations is needed only for the arithmetic
representation below.

Let $E$ be any set of initial valuations from which the pCFG is AST, and
let $\mathit{Inv}$ contain exactly the configurations reachable from $E$
by positive-probability paths. Every configuration in $\mathit{Inv}$ is
AST: otherwise, a positive-probability path to a configuration with
positive divergence probability would give positive divergence probability
from an initial state. Set $U(\gamma)$ to the length of a shortest
positive-probability path from $\gamma$ to the program exit. Such a finite
path exists by AST. Hence $U$ is zero exactly at the exit and strictly
decreases at the unique successor of a deterministic nonterminal
configuration and at some successor of a probabilistic configuration.
The minimum of $1$ and the positive probabilistic edge probabilities
of the finite pCFG gives a common positive progress bound.

It remains to construct a compatible $V$. Fix an injective natural-number
encoding of configurations, computable for finite rational state tuples. For $\gamma\in\mathit{Inv}$, let
$R(\gamma,n)$ be the probability of visiting a nonterminal configuration
with code at least $n$ before termination, including the starting
configuration. For terminal $\gamma$, set $R(\gamma,n)=0$.
For fixed $n$, this function is a nonnegative supermartingale: outside
the target set the first-step equation is an equality, and in the target
set its value is $1$. Moreover,
\[
 \lim_{n\to\infty}R(\gamma,n)=0.
\]
Indeed, an almost-surely terminating run visits only finitely many
configurations before its first exit, so the corresponding decreasing
tail events have an intersection of probability zero.

For every $k\geq1$, choose the least integer $n_k\geq k$ such that
\[
 R(\gamma,n_k)\leq2^{-k}
 \quad\text{for all }\gamma\in\mathit{Inv}
                     \text{ with }\operatorname{code}(\gamma)\leq k.
\]
Only finitely many configurations are constrained for each $k$, so this
integer exists. Define
\[
 V(\gamma)=[\gamma\text{ is nonterminal}]
              +\sum_{k=1}^{\infty}R(\gamma,n_k).
\]
For a fixed configuration with code $i$, the summands with $k\geq\max(1,i)$
are bounded by $2^{-k}$; thus $V$ is finite everywhere. The indicator
is itself a supermartingale because the exit is absorbing. Nonnegative
summation preserves the supermartingale inequalities, and the indicator
term makes $V$ positive at every nonterminal configuration and zero at
the exit. Given $r\geq0$, choose an integer $m>r$. Any nonterminal
configuration with code at least $\max\{n_1,\ldots,n_m\}$ belongs to all
first $m$ target sets, and hence has $V\geq1+m>r$. Consequently,
$V_{\leq r}$ contains only finitely many nonterminal configurations.
As $U=0$ at terminal configurations, $U$ is bounded on every sublevel
set. This proves semantic completeness for arbitrary $E$, including
infinite sets, without any uniform bound on initial certificate values.
The empty initial set is handled by the empty invariant.

\paragraph{Arithmetic representation.}
If $E$ is arithmetically definable, so is $\mathit{Inv}$: quantify over
an initial state in $E$ and an encoded finite feasible path. The same
encoding defines the graph of the shortest-path function $U$. For a finite horizon $t$, let $h_t(\gamma,n)$ be the probability of
hitting a nonterminal configuration of code at least $n$ by step $t$
before termination. It is a computable rational number, uniformly in
its arguments, because the transition tree to depth $t$ is finite.
Since $R(\gamma,n)=\sup_t h_t(\gamma,n)$, for rational $q$,
\begin{align*}
 q<R(\gamma,n)&\ \Longleftrightarrow\ \exists t:\ q<h_t(\gamma,n),\\
 R(\gamma,n)\leq q&\ \Longleftrightarrow\ \forall t:\ h_t(\gamma,n)\leq q.
\end{align*}
Thus the defining property and minimality of $n_k$ give an arithmetic
graph for the sequence, using the arithmetic predicate for
$\mathit{Inv}$. To encode the strict lower cut of $V$, write
$d_\gamma=[\gamma\text{ is nonterminal}]$. Then $q<V(\gamma)$ iff
there exist $m\in\mathbb N$ and an encoded rational sequence
$(q_1,\ldots,q_m)$ such that
\[
 q<d_\gamma+\sum_{k=1}^m q_k
 \quad\text{and}\quad
 q_k<R(\gamma,n_k)\text{ for every }1\leq k\leq m.
\]
The empty sequence is allowed. This equivalence follows from convergence
of the nonnegative series and density of the rationals for each finite
partial sum. It supplies one arithmetic formula, rather than a separate
formula for each bound or configuration. No effective enumeration of
$E$, and no computation of the infinite sum, is required. These are
expressibility statements relative to true arithmetic.
